\section{Exact morphisms, rigid charts and principal quantisations}
\label{sec:setup}

This section fixes conventions and records the inputs from the preceding
iteration.  The new globalization arguments begin in \cref{sec:globalization}.

\subsection{The exact morphism category}

Let $A$ be a basic finite-dimensional $\kk$-algebra with indecomposable
projective right modules $P_1,\ldots,P_n$.  We write
\[
 K:=K_0^{\mathrm{sp}}(\proj A)
   =\bigoplus_{i=1}^n\mathbb Z[P_i]
 \simeq\mathbb Z^n.
\]
The exact category
\[
 \MM=\Mor(\proj A)
\]
has objects
\[
 X=(P^-\xrightarrow fP^0)
\]
and degreewise split conflations.  Set
\[
 Z_P=(P\longrightarrow0),
 \qquad
 U_P=(0\longrightarrow P),
 \qquad
 K_P=(P\xrightarrow{1_P}P).
\]
The term functor induces
\begin{equation}
 K_0(\MM)\xrightarrow{\sim}L:=K\oplus K,
 \qquad
 [P^-\xrightarrow fP^0]\longmapsto([P^-],[P^0]).
 \label{eq:term-lattice}
\end{equation}
We use the ordered basis
\[
 (Z_1,\ldots,Z_n,U_1,\ldots,U_n).
\]
The split-difference map is
\begin{equation}
 D=(-I_n\ I_n):L\longrightarrow K,
 \qquad
 D(p,q)=q-p.
 \label{eq:D-map}
\end{equation}
Thus
\[
 D[X]=[P^0]-[P^-]=:I(X),
\]
the split projective index.  Its kernel is the diagonal lattice
\begin{equation}
 C:=\ker D=\{(r,r):r\in K\},
 \label{eq:contractible-lattice}
\end{equation}
generated by the contractible classes $[K_{P_i}]$.

For projective morphisms
$X=(P^-\xrightarrow fP^0)$ and
$Y=(Q^-\xrightarrow gQ^0)$, the exact Hom--Ext complex gives
\begin{equation}
\begin{aligned}
0\longrightarrow\Hom_{\MM}(X,Y)
&\longrightarrow
 \Hom(P^-,Q^-)\oplus\Hom(P^0,Q^0)\\
&\xrightarrow{(a,b)\mapsto ga-bf}
 \Hom(P^-,Q^0)
 \longrightarrow\Ext^1_{\MM}(X,Y)\longrightarrow0.
\end{aligned}
\label{eq:hom-ext-complex}
\end{equation}
Consequently the exact Euler form depends only on the two projective terms.
If
\[
 H_{ij}=\dim_\kk\Hom_A(P_i,P_j),
\]
then its matrix and its skew-symmetrization in the term basis are
\begin{equation}
 E_{\MM}=
 \begin{pmatrix}H&-H\\0&H\end{pmatrix},
 \qquad
 \Lambda_{\MM}=E_{\MM}^T-E_{\MM}
 =\begin{pmatrix}
 H^T-H&H\\-H^T&H^T-H
 \end{pmatrix}.
 \label{eq:exact-euler-skew}
\end{equation}

Over $\mathbb F_q$, $q=v^2$, let $\delta_X$ denote the
extension-normalized Hall basis and put
\[
 \eps_X=v^{-\langle X,X\rangle_{\MM}}\delta_X.
\]
If $\Ext^1_{\MM}(X,Y)=0$, then
\begin{equation}
 \eps_X\diamond\eps_Y
 =v^{\Lambda_{\MM}([X],[Y])}\eps_{X\oplus Y}.
 \label{eq:split-Hall}
\end{equation}
For any integral skew form $\Omega$ on $L$, the bicharacter twist
\begin{equation}
 a\star_\Omega b
 =v^{\Omega(\deg a,\deg b)}a\diamond b
 \label{eq:Hall-twist}
\end{equation}
is associative.  The skew-normalized exact integration is an algebra map to
the quantum torus whose form is $\Lambda_{\MM}+\Omega$.  These statements use
only bilinearity and the standard exact integration theorem
\cite{ZhangFixedSize}.

\subsection{Homology and reduced morphisms}

For completeness we recall the homological normal form.  Let
\[
 X=(P^-\xrightarrow fP^0),
 \qquad N=H^0(X)=\coker f.
\]
Write
\[
 C_N=(P_1^N\longrightarrow P_0^N)
\]
for a minimal projective presentation of $N$.  There are uniquely determined
projectives $P_X,Q_X$, up to isomorphism, such that
\begin{equation}
 X\simeq C_N\oplus Z_{P_X}\oplus K_{Q_X},
 \qquad
 H^{-1}(X)\simeq\Omega_A^2N\oplus P_X.
 \label{eq:homology-normal-form}
\end{equation}
The shifted-projective residue is obtained by actual Krull--Schmidt
cancellation,
\begin{equation}
 P_X=H^{-1}(X)\ominus\Omega_A^2H^0(X),
 \label{eq:residual-cancellation}
\end{equation}
and
\begin{equation}
 I(X)=g_A(N)-[P_X].
 \label{eq:index-homology}
\end{equation}
The symbol $\ominus$ in \cref{eq:residual-cancellation} is not a formal
Grothendieck difference: it subtracts multiplicities of indecomposable
Krull--Schmidt summands.  In particular, taking the largest projective direct
summand of $H^{-1}(X)$ is not valid when $\Omega_A^2N$ itself has projective
summands.

The reduced object
\[
 X_{\mathrm{red}}=C_N\oplus Z_{P_X}
\]
is therefore reconstructed by the full homology pair.  Li's correspondence
identifies reduced rigid objects with support $\tau$-rigid pairs and maximal
rigid reduced objects with support $\tau$-tilting pairs \cite{Li,AIR}.

\subsection{Cluster-tilted realization and two lattice labels}

Assume from now on that
\[
 A\simeq\End_{\CC}(T)
\]
for a basic cluster-tilting object $T=T_1\oplus\cdots\oplus T_n$ in a
Hom-finite $2$-Calabi--Yau cluster category, and fix the quantum seed attached
to $T$.  We use right $A$-modules and the sign convention
\[
 C_T=-B_0,
\]
where $B_0$ is the actual initial exchange matrix in the natural morphism
seed.  The two-term object represented by $X$ is
\[
 \Psi_T(X)=\Sigma^{-1}\Cone(f).
\]
For a reduced rigid morphism, its cluster $g$-vector is reconstructed from
homology by
\begin{equation}
 g_T(X)=C_T\dimv H^0(X)-I(X).
 \label{eq:cluster-g-homology}
\end{equation}
Its classical morphism character is
\begin{equation}
 CC_T^{\mathrm{mor}}(X)
 =\sum_{\mathbf e}
   \chi\!\left(\Gr_{\mathbf e}H^0(X)\right)
   x^{C_T(\dimv H^0(X)-\mathbf e)-I(X)}.
 \label{eq:morphism-character}
\end{equation}
Thus the split index $I(X)$ and the cluster label $g_T(X)$ are different
integer vectors attached to the same object.

Let
\[
 R=f_1\oplus\cdots\oplus f_n
\]
be a reachable ordered basic reduced maximal rigid morphism object.  Define
\begin{equation}
 \Delta_R=(I(f_1)\ \cdots\ I(f_n)),
 \qquad
 G_R=(g_T(f_1)\ \cdots\ g_T(f_n)).
 \label{eq:Delta-G}
\end{equation}
Both matrices lie in $\mathrm{GL}_n(\mathbb Z)$.  The columns of $\Delta_R$
are the $g$-vectors of the suspended companion cluster-tilting object, while
the columns of $G_R$ are those of the original cluster-tilting object.

The presentation-to-cluster transfer in the chamber $R$ is
\begin{equation}
 \mathsf L_R:=G_R\Delta_R^{-1}\in\mathrm{GL}_n(\mathbb Z).
 \label{eq:chamber-transfer}
\end{equation}
If $\delta=\Delta_Rm$ with $m\in\mathbb N^n$, then the unique rigid open
orbit in the corresponding presentation space is represented by
\[
 F_{R,m}=\bigoplus_{i=1}^nf_i^{m_i},
\]
and
\begin{equation}
 \Gamma_T(\delta)=G_Rm=\mathsf L_R\delta.
 \label{eq:fan-transport}
\end{equation}
The local maps agree on common faces.  Hence they define a continuous
piecewise-integral-unimodular map on the reachable presentation fan.  The
fan--orbit equivalence identifies its domain as
\begin{equation}
 \delta\in\lvert\Fan_\tau(A)\rvert
 \quad\Longleftrightarrow\quad
 \Pres_A(\delta)\text{ has a rigid open orbit}
 \quad\Longleftrightarrow\quad
 e_A(\delta,\delta)=0.
 \label{eq:fan-orbit}
\end{equation}
Here $e_A$ is the generic Derksen--Fei $E$-invariant
\cite{DerksenFei,DIJ}.

\subsection{Principal-compatible quantum forms}

At the initial seed put
\[
 \widetilde B_0=\binom{B_0}{I_n}.
\]
For $d>0$, the canonical principal form at a chart $R$ is
\begin{equation}
 \widetilde\Lambda_R^{\mathrm{can}}
 =d\begin{pmatrix}
      0&-G_R^T\\
      G_R&B_0^T
    \end{pmatrix}.
 \label{eq:canonical-principal-form}
\end{equation}
It satisfies
\[
 \widetilde B_R^T\widetilde\Lambda_R^{\mathrm{can}}
 =d(I_n\ 0),
 \qquad
 \widetilde B_R=\binom{B_R}{C_R}.
\]
Every principal-compatible integral skew form with the same diagonal
compatibility can be transported from an initial parameter
$S\in\Skew_n(\mathbb Z)$ as
\begin{equation}
 \widetilde\Lambda_R(S)
 =\widetilde\Lambda_R^{\mathrm{can}}
  +\binom{G_R^T}{-B_0^T}
    S
    \begin{pmatrix}G_R&-B_0\end{pmatrix}.
 \label{eq:principal-affine-family}
\end{equation}
Its mutable--mutable block is
\begin{equation}
 \Lambda_R^{\mathrm{mut}}(S)=G_R^TSG_R.
 \label{eq:mutable-block}
\end{equation}
Formula \cref{eq:principal-affine-family} is the homology-adapted version of
the complete principal quantisation space; compare
\cite{GellertLampe}.  It is valid whether or not $B_0$ is invertible.

For each column $\delta$ of $\Delta_R$, define its reduced term lift
\[
 \xi(\delta)=\binom{(-\delta)_+}{\delta_+}\in\mathbb N^{2n},
\]
and put
\begin{equation}
 \Xi_R=(\xi(I(f_1))\ \cdots\ \xi(I(f_n))).
 \label{eq:term-lift}
\end{equation}
Then
\begin{equation}
 D\Xi_R=\Delta_R.
 \label{eq:D-Xi}
\end{equation}
The problem addressed in the next section is to decide when the chamber forms
\cref{eq:mutable-block} are the restrictions of one skew form on the exact
term lattice $L$.
